\section{现代 filtering：Dolma、DCLM、Nemotron-CC}

Raw web text 规模巨大，但有效训练预算有限，因此本节比较现代 pipeline 如何把规则、轻量 classifier、强模型打分与 dedup 组合成过滤级联。关键不只是最终保留率，而是正负样本如何定义“质量”、不同 domain 的 false positive/negative，以及过滤成本能否覆盖万亿 token 规模。

\begin{figure}[H]
\centering
\includegraphics[width=0.82\textwidth]{dolma-overview.png}
\caption{Dolma 数据概览：Reddit、PeS2o、C4、Project Gutenberg、Wikipedia/Wikibooks 等组成。}
\figsource{Dolma overview image，本地化后供 CS336 Lecture 13 使用。}
\end{figure}

\begin{knowledgebox}{讲义提醒：过滤器也需要自己的 held-out audit set}
不能只用训练过滤器的正负例评估它。应按语言、域名、文档长度、来源类型与风险类别分层抽样，人工检查被保留和被删除的两侧；否则 classifier 很容易把排版、网站模板或 benchmark 风格当成内容质量。这是由课程展示的多种 filtering recipe 推导出的工程审计实践。
\end{knowledgebox}

\begin{knowledgebox}{读图：Dolma 的 pipeline}
Dolma 的 Common Crawl 处理包括语言识别、Gopher/C4 规则过滤、toxicity filtering、Bloom filter deduplication。它体现了现代开源数据集的折中：不用强模型过滤以减少 bias，但仍大量依赖启发式。
\end{knowledgebox}

\begin{figure}[H]
\centering
\includegraphics[width=0.9\textwidth]{dclm-filter.png}
\caption{DCLM：从 240T token DCLM-pool 中用 quality classifier 过滤出 baseline 数据。}
\figsource{CS336 Lecture 13 官方源引用图。}
\end{figure}

\begin{figure}[H]
\centering
\includegraphics[width=0.78\textwidth]{dclm-quality.png}
\caption{DCLM quality classifier 与其他 filtering 方法的比较。}
\figsource{CS336 Lecture 13 官方源引用图。}
\end{figure}

\begin{knowledgebox}{读图：model-based filtering 的优缺点}
DCLM 使用正例 OpenHermes/ELI5、负例 RefinedWeb 训练 fastText classifier，在大池子上过滤。优点是质量提升明显；风险是 classifier 会继承正负例选择偏差，把“像某些高质量数据”误当成“真正有用数据”。
\end{knowledgebox}

\begin{importantbox}{Model-based filtering 的机制}
把一个小的人工/模型选定数据集当作“好文本”正例，把普通 web 文本当作负例，训练 classifier 后打分整个 Common Crawl pool。这个过程实质上是在定义“模型应该多学哪类文本”。因此 classifier 的训练集就是隐含的数据价值观。
\end{importantbox}

\begin{knowledgebox}{课堂提示：正负例已经在定义“质量”}
DCLM 的正例包括 OpenHermes-2.5 和 ELI5，负例来自 RefinedWeb。老师把样本来源直接列出来，是为了让读者看到：classifier 并非从真空中发现质量，它学习的是“更像 instruction/解释型文本，还是更像普通 web 文本”。读结果时应同时审计正负例构成，而不能只比较最终 benchmark 分数。
\end{knowledgebox}

\begin{figure}[H]
\centering
\includegraphics[width=0.9\textwidth]{nemotron-results.png}
\caption{Nemotron-CC：在不过度过滤的前提下保留更多 token，并结合 classifier ensemble 和 synthetic rephrasing。}
\figsource{CS336 Lecture 13 官方源引用图。}
\end{figure}

\begin{knowledgebox}{读图：Nemotron-CC 的动机}
FineWebEdu 和 DCLM 可能移除 90\% 数据，质量高但 token 少。Nemotron-CC 试图“need more tokens but preserve quality”：用 Nemotron-340B-instruct 打分并蒸馏 classifier，结合 DCLM classifier，对低质量数据重写，对高质量数据生成任务。
\end{knowledgebox}

\begin{knowledgebox}{合成数据和重写的风险}
用 LM rephrase 低质量数据或从高质量数据生成 QA pairs，可以提高可学信号密度；但也可能引入 teacher model 的风格偏差、事实错误、重复模板和版权继承问题。合成数据应记录生成模型、prompt、过滤器和去重策略。
\end{knowledgebox}

